\section{Tensor prices and wire distances}\label{sec:prices}

A tensor network supplies a coefficient array by sending labels along
wires and contracting tensors at their endpoints. A price assigns a
number to the array once its indices have been placed at sites. Local
linear maps can only decrease that number, and the prices of independent
wires add. Thus the wires in a network must together pay at least the
price of its output. We will also need the cost of two operations on a
network, moving an output index and summing several possible supplies.
Throughout this section and Part~\ref{part:no}, we use natural logarithms
for prices. This multiplies the base-two prices of \cite{oaG} by
$\log 2$ and preserves the units of wire distance.

\subsection{Local maps, placements and limits of tensors}

Fix a field $K$ and a finite set $J$ of at least two sites. A
\emph{restriction} acts on a tensor by a linear map at each site
(or by precomposition when the tensor is viewed as a multilinear
form). Two operations combine tensors on $J$. The tensor product
groups the corresponding factors at each site. The direct sum places
the two tensors in disjoint blocks of the local spaces. Up to mutual
restriction, these operations form the tensor semiring in Strassen's
asymptotic spectrum \cite{Strassen}.

A \emph{spectral point} $\varphi$ is a nonnegative real-valued
function on this semiring. It is monotone under restriction,
multiplies on tensor products, adds on direct sums, and preserves
zero and one. Matrix rank after flattening across a fixed nontrivial
bipartition is one example. Choose a scale $D>0$ and put
$H(T)=D\log\varphi(T)$, with $H(0)=-\infty$. Every nonzero tensor
restricts to a nonzero simple tensor, so $H(T)\ge0$ for $T\ne0$.

To share one label with $q$ possible values among the sites of a
nonempty $A\subseteq J$, use the diagonal tensor
\[
 \Delta_A(q)=\sum_{i=1}^q\bigotimes_{a\in A}e_i,\qquad
 \delta(A)=H(\Delta_A(2))/\log 2 ,
\]
with one-dimensional slots outside $A$. Tensor products multiply
label sizes, and restriction can reduce them. Comparing powers of
$q$ with powers of two therefore gives $H(\Delta_A(q))=\delta(A)\log q$.
In particular $\delta(J)=D$, since $\Delta_J(q)$ is the direct sum of
$q$ unit tensors. For $a\ne b$, the two-site diagonal is a wire from
$a$ to $b$. Its rate $d(a,b)=\delta(\{a,b\})$ defines a pseudometric
when $d(a,a)=0$. Joining two wires gives the triangle inequality.
A diagonal restricts to a wire between any pair of its sites.
Conversely, equating the labels on the edges of a tree spanning $A$
supplies the diagonal, even if the tree uses auxiliary vertices.
Consequently
\[
    \operatorname{diam}(A)\leq\delta(A)\leq
    \sum_{uv\ \mathrm{in\ the\ tree}}d(u,v).
\]

In a \emph{spectral space}, every finite set $A$ of positive diameter
carries a price $H_A=\delta(A)\log\varphi_A$ that induces the given
wire distances. The prices agree when one adds one-dimensional slots,
and a nonzero tensor supported at a single site has price zero. The
metric quotient of the sites of any spectral point carries this
structure \cite[Section~\ext{G}{sec:spectra}]{oaG}. A \emph{placement}
$x=(x_i)$ assigns each index of a tensor $T$ to a site. Group indices
assigned to the same site and write $H(T;x)$ for the resulting price.
A local map may act jointly on all indices at its site.

\par\Needspace{6\baselineskip}
\begin{lemma}[Price inequalities]\label{price:rules}
Let $T\ne0$ have local dimensions $q_i$ and let $x,y$ be two placements
in one spectral space. Then
\begin{equation}\label{price:transport}
   |H(T;x)-H(T;y)|\leq\sum_i d(x_i,y_i)\log q_i .
\end{equation}
For $L\ge1$ tensors of the same format at sites in $A$,
\begin{equation}\label{price:sum}
    H\left(\sum_{j=1}^L T_j\right)
       \leq\max_j H(T_j)+\delta(A)\log L.
\end{equation}
At fixed placement and format over an algebraically closed field, every
price sublevel is Zariski closed.
\end{lemma}

\begin{proof}
A wire of dimension $q_i$ can transport index $i$ from $x_i$ to $y_i$
by a contraction. Reversing the move proves Equation~\eqref{price:transport}.
Folding the local summand spaces takes a direct sum to an ordinary
sum. Apply additivity of $\varphi_A$ and take logarithms to obtain
Equation~\eqref{price:sum}. At a single site every nonzero sum has
price zero.

For closedness, let $U$ lie in the closure of a set of tensors of
price at most $b$ in an ambient tensor space of dimension $v$.
For every $h\ge1$, the map $T\mapsto T^{\otimes h}$ is polynomial.
The span of its values on the set is closed and hence contains
$U^{\otimes h}$. Expanding in degree-$h$ monomials in the entries
of $T$ bounds the dimension of this span by
$\binom{v+h-1}{v-1}$. Express $U^{\otimes h}$ in a basis of such
powers. Scaling a tensor by a nonzero constant preserves its price.
Equation~\eqref{price:sum} thus gives
\[
 hH(U)\le hb+\delta(A)\log\binom{v+h-1}{v-1},
\]
where $A$ is the set of occupied sites. Dividing by $h$ and letting
$h\to\infty$ proves the claim
\cite[Lemma~\ext{G}{lem:spectral-closure}]{oaG}
(see also \cite{SpectralSublevels}).
\end{proof}

For a network with fixed output sites, applying its local coefficient
tables to the product of its wire tensors supplies the output by
restriction. The sum of the wire prices therefore bounds the output
price, including in a limit at the same format. If a network shares
one label across several sites, its contribution is the corresponding
diagonal price. Equation~\eqref{price:sum} gives the same charge for
the common label selecting a summand.

\subsection{The spaces of wire distances}

A Hadamard space is a complete geodesic $\operatorname{CAT}(0)$ space.
Every spectral space extends isometrically to a Hadamard spectral
space that preserves its prices. The price of a nonzero tensor is
convex as its indices move simultaneously along constant-speed
geodesics. Three operations preserve spectral realizability. One can
rescale the metric and all prices by a common positive factor, pass
to a closed convex subspace, or take a pointed ultralimit
\cite[Theorem~\ext{G}{thm:geometry}, Proposition~\ext{G}{prop:joint-convexity}
and Section~\ext{G}{sec:limits}]{oaG}. For an ultralimit,
Equation~\eqref{price:transport} bounds the price on bounded
configurations and ensures independence of the representing sequences.

Let $\mathcal X_\chi$ denote the class of Hadamard spaces admitting
spectral prices over a field of characteristic $\chi$. The geometry
depends only on $\chi$. Indeed, scalar extension preserves the
possible simultaneous spectral values of tensors over the subfield
\cite[Lemma~\ext{G}{lem:comparison}]{oaG}, and every binary wire is
defined over the prime field. Finite wire-distance tables suffice to
recover a spectral structure on a Hadamard space. Fix a basepoint
and take a pointed ultralimit of realizations of its finite subsets
containing that point. An ultrafilter containing every tail gives an
isometric copy of the space, to which the limiting prices restrict.
We may therefore work over an uncountable algebraically closed field
of the chosen characteristic.

A class of metric spaces \emph{admits ports} if regular multigraphs
of a single positive degree and unbounded orders satisfy a uniform
Poincar\'e inequality for squared distances in all its spaces, as in
Section~\ref{sec:intro}. We allow symmetric nonnegative integral arc
counts of the given row sum, including diagonal entries. The condition
for $\mathcal X_\chi$ is precisely port admission by $\chi$
\cite[Definition~\ext{G}{def:ports}]{oaG}. Taking a metric quotient,
duplicating sites, or rescaling their distances preserves the inequality.

\subsection{The cost of sharing the summand index}

The balanced-cut decomposition in Part~\ref{part:no} will have at most
$2^d$ summands for $d$ row sites in a unit ball. By
Equation~\eqref{price:sum}, the charge for choosing the summand
depends on the diagonal rate of the row sites. The tree bound can grow
linearly in the number of sites. In the spaces we use, the following
estimate reduces this growth to the square root of its logarithm.

If $\chi$ has no ports, Theorem~\ref{prod:main} gives a subclass of
$\mathcal X_\chi$ that still has no ports and is closed under finite
$\ell_2$ products, positive rescaling, closed convex subspaces and
pointed ultralimits. This subclass contains a nontrivial segment,
since a class of one-point spaces admits every port family. Rescaling
segments and taking a pointed limit gives a line. Products of lines
and pointed limits then give a copy of the separable Hilbert space
$\ell_2$ (pass to the closed convex subspace generated by the
coordinate directions). Thus for $Y$ in the subclass we can place a
configuration in $Y\times\ell_2$ at Hilbert coordinate zero and choose
spectral prices on the product.

When $Y$ is a metric cone with tip $o$, these prices can be made
positively homogeneous under radial dilations. For $T\ne0$ replace
$H(T;x)$ by
\[
       \lim_{s\downarrow0} H(T;s x)/s .
\]
Joint convexity and the zero price at the tip give the limit. It is
also an ultralimit of the rescaled spectral prices, so it has the
same wire metric and satisfies the spectral axioms. This construction
applies to products of cones as well.

\begin{lemma}[Diagonal price in a ball]\label{price:diagonal}
Let $A$ consist of $n\ge2$ sites in a ball of radius $r$ in a geodesic
metric space $Y$. Equip $Y\times\ell_2$ with spectral prices, and
place the sites at Hilbert coordinate zero. Then
\[
  \delta(A)\leq \sqrt{2\pi}\,r n\,
  \phi\bigl(\Phi^{-1}(1-1/n)\bigr)=O(r\sqrt{\log n}),
\]
where $\Phi$ and $\phi$ are the standard normal distribution function
and density.
\end{lemma}

\begin{proof}
Rescale to $r=1$, since the claim at $r=0$ is immediate. Write $o$
for the ball center. We will compare the marginal cost of adding
each $a\in A$ to a ball in the Hilbert coordinates through $o$. Each
comparison takes place in a Euclidean patch, the product of a single
geodesic segment $[o,a]$ with a finite-dimensional Euclidean space.

First we show that on any finite-dimensional convex Euclidean patch
the diagonal rate of a finite nonempty set $V$ equals its mean width,
\[
 W(V)=\sqrt{2\pi}\,\EE\max_{v\in V}\langle G,v\rangle,\qquad
 G\ \text{a standard Gaussian vector}.
\]
The flat law of \cite[Theorem~\ext{G}{thm:flatlaw}]{oaG} gives this
identity on an entire Euclidean space. At a relative interior point
$y$ of a patch, each tangent rescaling limit contains its whole
translation space. Hence for every finite nonempty configuration $E$
of directions in that space,
\[
          \lim_{h\downarrow0}\delta(y+hE)/h=W(E).
\]
To pass from tangent configurations to $V$, work in its convex hull.
Adjoining any point of this hull preserves price. Indeed, the added
leg can be copied at any existing site, and its price is convex
in its position. Discarding the added leg gives the reverse inequality.
Choose a relative interior point as origin and
adjoin it to $V$. Thus $t\mapsto\delta(tV)$ has right derivative
$W(V)$ at $t=0$.

To bound its derivative at $0<t<1$, adjoin $(t+h)v$ to $tV$ for
each $v\in V$ in turn. Write $y=tv$ and $E=tV-y$, so $0\in E$.
For fixed $M>0$ and small $h>0$, the smaller configuration
$y+MhE$ lies in the convex hull of $tV$ and contains $y$. Adjoin
it at zero additional cost. Submodularity on intersecting sets
\cite[Proposition~\ext{G}{lem:submod}]{oaG} now bounds the increment
from adding $y+hv$ by the same increment starting from $y+MhE$
alone. The point $y$ is interior relative to the hull of $V$. After
division by $h$, the latter increment tends to
\[
             W(ME\cup\{v\})-W(ME).
\]
Let $M\to\infty$. Since $0\in E$, dominated convergence gives the bound
\[
 \sqrt{2\pi}\,\EE\!\left[
   \langle G,v\rangle_+\,
   \mathbf1_{\{\max_{e\in E}\langle G,e\rangle=0\}}\right].
\]
The event means precisely that $v$ maximizes $\langle G,\cdot\rangle$
on $V$. Almost surely the maximizing point is unique, and its value
is nonnegative since $0$ is in the hull. Summing these increments over
$v$ bounds the right derivative of $\delta(tV)$ by $W(V)$.
Convexity and continuity now give $\delta(V)=W(V)$.

For $t\ge0$, take a finite net $F$ in the Euclidean ball of radius
$t/\sqrt j$ in $j$ Hilbert directions through $(o,0)$. Monotonicity
and submodularity give
\[
 \delta(A)\leq\delta(F)+
        \sum_{a\in A}\bigl(\delta(F\cup\{a\})-\delta(F)\bigr).
\]
Use the patch formula for $F$ and $F\cup\{a\}$. Sending the mesh to
zero at fixed $j$ gives the Gaussian maximum $t\|G_j\|/\sqrt j$ on
the ball, where $G_j$ is a standard $j$-dimensional Gaussian. The same Gaussian linear functional has
value at $a$ equal to $d(o,a)G$ for an independent standard normal $G$.
The ball maximum tends in mean square to $t$ as $j\to\infty$, so
\[
 \delta(A)\leq\sqrt{2\pi}\left(t+
          \sum_{a\in A}\mathbb E(d(o,a)G-t)_+\right)
       \leq\sqrt{2\pi}\left(t+n\mathbb E(G-t)_+\right).
\]
Taking $t=\Phi^{-1}(1-1/n)$ and integrating the normal density gives
the stated expression. Taking $t=\sqrt{2\log n}$ gives the order bound.
\end{proof}

The leading constant is sharp even on a Hilbert space. For $n$
orthogonal unit vectors the flat law gives
$\delta(A)=\sqrt{2\pi}\,\EE\max_{i\le n}G_i$ for independent standard
normals $G_i$. The normal tail estimate gives
$\EE\max_iG_i\sim\sqrt{2\log n}$, so both sides of the bound have
leading term $2\sqrt{\pi\log n}$ at radius one.
